% Layout gallery only: all numbers and research examples below are synthetic.
\documentclass[aspectratio=169,10pt]{beamer}
% Edit this file before loading paper-slides.sty.
\newif\ifChineseSlides
\ChineseSlidesfalse % Set true AND compile with XeLaTeX for Chinese slides.
\newcommand{\LatinFontPackage}{FiraSans} % Empty = TeX default; sans font package.
\newcommand{\CJKFontSet}{fandol} % TeX-distributed fonts, used by ctex.

\definecolor{SlideBlue}{HTML}{00498A}
\definecolor{SlideBlueDark}{HTML}{1F3A5F}
\definecolor{SlideRed}{HTML}{C8102E}
\definecolor{SlideLight}{HTML}{CCDEEC}
\definecolor{SlideGreen}{HTML}{2E7D32}
\definecolor{SlideGray}{HTML}{595959}
\definecolor{SlideInk}{HTML}{1A1A1A}

\title[Research Talk]{A Clear Research Story}
\subtitle{Reusable academic presentation layouts}
\author{Anonymous Authors}
\newcommand{\Presenter}{Anonymous Presenter}
\institute{} % Full institution(s); respect anonymous-review rules.
\newcommand{\HeaderInstitution}{} % Optional short institution label.
\newcommand{\Contact}{} % Optional public contact; blank by default.
\newcommand{\TalkType}{Research Presentation}
\date{}

\usepackage{paper-slides}

\begin{document}
% Layout 1: title. The title page counts as PDF page / frame 1.
\begin{frame}[plain]
  \titlepage
\end{frame}

% Layout 2: two-column explanation.
% Source: synthetic design example, not a claim about a real paper.
\begin{frame}{Start with the research question}
  \takeaway{A clear question connects the motivation to the proposed idea.}
  \begin{columns}[t,onlytextwidth]
    \begin{column}{0.47\textwidth}
      \kw{What is difficult?}
      \begin{itemize}
        \item Describe the concrete setting.
        \item Show why the current approach fails.
        \item State the condition that matters.
      \end{itemize}
    \end{column}
    \begin{column}{0.47\textwidth}
      \kw{What changes?}
      \begin{itemize}
        \item Introduce one central idea.
        \item Explain its intended benefit.
        \item Name the evidence needed to test it.
      \end{itemize}
    \end{column}
  \end{columns}
  \vspace{0.8em}
  \begin{block}{Research question}
    Can the proposed change improve the outcome under the same evaluation conditions?
  \end{block}
\end{frame}

% Layout 3: supplemental intuition; use a core original for the main method overview.
% A simplified redraw does not replace an available, useful original for convenience.
\begin{frame}{Show the idea before the details}
  \takeaway{Explain the input, the key operation, and the output.}
  \begin{center}
    \begin{tikzpicture}[>=Stealth,node distance=1.0cm,
      box/.style={draw=SlideBlue,fill=SlideLight!45,rounded corners,
        minimum width=2.8cm,minimum height=1.1cm,align=center}]
      \node[box] (input) {Input\\and assumptions};
      \node[box,right=of input] (method) {Key\\operation};
      \node[box,right=of method] (output) {Output\\and evaluation};
      \draw[->,thick,SlideBlue] (input)--(method);
      \draw[->,thick,SlideBlue] (method)--(output);
    \end{tikzpicture}
  \end{center}
  \vspace{0.8em}
  \begin{itemize}
    \item Name the assumption that makes the operation meaningful.
    \item Use the next page to explain its essential mechanism.
  \end{itemize}
\end{frame}

% Layout 4: equation and algorithm intuition; synthetic optimization example.
\begin{frame}{Explain the essential relationship}
  \takeaway{Define the symbols and conditions before interpreting the equation.}
  \begin{columns}[t,onlytextwidth]
    \begin{column}{0.48\textwidth}
      \kw{Objective}
      \[
        \theta^* = \arg\min_{\theta\in\Theta}
        \frac{1}{n}\sum_{i=1}^{n}\ell(f_\theta(x_i),y_i)
      \]
      \begin{itemize}
        \item $n$: number of training examples.
        \item $\ell$: per-example loss.
        \item $\Theta$: feasible parameter set.
      \end{itemize}
    \end{column}
    \begin{column}{0.46\textwidth}
      \begin{block}{Procedure sketch}
        \begin{enumerate}
          \item Evaluate the current solution.
          \item Update under the constraints.
          \item Stop at the stated criterion.
        \end{enumerate}
      \end{block}
      An objective alone does not guarantee that an algorithm finds its global optimum.
    \end{column}
  \end{columns}
\end{frame}

% Layout 5: readable table with honest comparison wording.
% Source: synthetic numbers. Derived: 55 - 50 = 5 percentage points.
\begin{frame}{Report gains and ties precisely}
  \takeaway{Synthetic example: +5 percentage points on A; a tie on B.}
  \begin{center}
    \renewcommand{\arraystretch}{1.3}
    \begin{tabular}{lrrl}
      \toprule
      Task & Baseline (\%) & Method (\%) & Difference \\
      \midrule
      A & 50.0 & \good{55.0} & +5.0 pp \\
      B & 40.0 & 40.0 & Tied \\
      \bottomrule
    \end{tabular}
  \end{center}
  \vspace{0.5em}
  All values here are \hot{illustrative}, not experimental results.
  \begin{itemize}
    \item State the model, dataset split, budget, and metric direction.
    \item A relative gain of 10\% on A differs from a gain of 5 pp.
  \end{itemize}
\end{frame}

% Layout 6: plot + interpretation. All values are synthetic.
\begin{frame}{Connect the figure to one observation}
  \takeaway{Synthetic example: more budget helps here, with smaller later gains.}
  \begin{columns}[c,onlytextwidth]
    \begin{column}{0.52\textwidth}
      \centering
      \begin{tikzpicture}[x=1cm,y=0.6cm,>=Stealth]
        \draw[->] (0,0)--(5.0,0) node[below right]{Budget};
        \draw[->] (0,0)--(0,4.7) node[above]{Score};
        \foreach \x in {1,2,3,4} \draw (\x,0.08)--(\x,-0.08) node[below]{\x};
        \foreach \y in {1,2,3,4} \draw (0.08,\y)--(-0.08,\y) node[left]{\y};
        \draw[very thick,SlideBlue] (1,1.5)--(2,2.7)--(3,3.2)--(4,3.4);
        \foreach \x/\y in {1/1.5,2/2.7,3/3.2,4/3.4}
          \fill[SlideBlue] (\x,\y) circle (2pt);
      \end{tikzpicture}
    \end{column}
    \begin{column}{0.43\textwidth}
      \kw{Interpretation}
      \begin{itemize}
        \item Point to the changing slope.
        \item Explain the cost of extra budget.
        \item Limit the conclusion to the observed range.
      \end{itemize}
    \end{column}
  \end{columns}
\end{frame}

% Layout 7: summary and limits.
\begin{frame}{Close the research story}
  \takeaway{Restate the contribution, the supporting evidence, and its boundary.}
  \begin{enumerate}
    \item \kw{Idea:} the one change the audience should remember.
    \item \kw{Evidence:} the observation that supports the claim.
    \item \kw{Boundary:} the conditions under which the claim holds.
  \end{enumerate}
  \vspace{0.7em}
  \begin{block}{Discussion}
    What additional evidence would make this conclusion stronger?
  \end{block}
\end{frame}

% Layout 8: end page; keep contact optional.
\begin{frame}{Discussion}
  \takeaway{Return to the central question.}
  \begin{center}
    {\Huge\color{SlideBlue}Thank you}\par\vspace{1em}
    {\large Questions and discussion}\par\vspace{0.7em}
    {\small\color{SlideGray}\Contact}
  \end{center}
\end{frame}


% Layout 9: original-figure overview. The figure leads; theme stays outside it.
% Source: self-created synthetic-method.tex, not a real paper figure.
% Treatment: original synthetic PDF, unchanged palette, arrows, and notation.
\begin{frame}{Start with the complete original figure}
  \takeaway{Use the full view to connect the proposal, check, and retained output.}
  \begin{center}
    \includegraphics[width=0.96\textwidth,height=4.5cm,keepaspectratio]{synthetic-method.pdf}
  \end{center}
  \small Read (a) $\rightarrow$ (b) $\rightarrow$ (c); follow the dashed feedback path.
  Panel (b) is enlarged next.
  \par\medskip
  {\footnotesize\color{SlideGray}Source: self-created synthetic example, panels (a)--(c).
  Original vector asset; no measured results.}
\end{frame}

% Layout 10: detail crop with explanation. Keep panel ID, legend, and context.
% Source: same synthetic-method.pdf, top-left crop [106, 0, 312, 136] pt.
% Treatment: crop only; TeX trim uses left/bottom/right/top in PDF points (bp).
% Incoming candidate comes from (a); outgoing yes arrow leads to (c).
\begin{frame}{Explain a focused detail in context}
  \takeaway{Panel (b): evaluation precedes acceptance; rejection returns to proposal.}
  \begin{columns}[c,onlytextwidth]
    \begin{column}{0.50\textwidth}
      \centering
      \includegraphics[width=\linewidth,height=4.7cm,keepaspectratio,
        trim=106bp 32bp 108bp 0bp,clip]{synthetic-method.pdf}
    \end{column}
    \begin{column}{0.46\textwidth}
      \kw{Locate panel (b) in the full view}
      \begin{enumerate}
        \item Candidate $h$ enters from panel (a).
        \item Evaluation supplies the score to the acceptance decision.
        \item The yes arrow leads to (c); the dashed path returns to (a).
      \end{enumerate}
      \small Legend: solid purple = forward; dashed orange = feedback.
    \end{column}
  \end{columns}
  \medskip
  {\footnotesize\color{SlideGray}Source: self-created synthetic example, panel (b).
  Cropped from the preceding full view; explanation and legend repeated outside.}
\end{frame}

% Add a references frame only when needed; replace the synthetic BibTeX entry.
% \begin{frame}{References}
%   \bibliographystyle{plain}
%   \bibliography{strings,refs}
% \end{frame}
\end{document}
